Population aging has become a major social issue in China, creating a significant digital divide as elderly users struggle with complex mobile terminal interfaces. This paper designs and implements SimpleDesktop, an age-friendly intelligent desktop system based on a micro-kernel plugin architecture and Android accessibility assistive technology.

The system adopts a three-layer micro-kernel architecture—host layer, bus layer, and plugin layer—decoupling business functions into independently loadable plugin units. A high-contrast, large-font interface is built with Jetpack Compose, supporting real-time font size and theme switching. A seven-stage state machine is designed using the Android Accessibility Service, combined with a dual-channel node retrieval algorithm and a focus anti-conflict mechanism, reducing WeChat video calling to a single tap. A multi-layer WiFi network guardian system ensures communication continuity. Voice interaction is achieved via the Android TTS engine with keyword matching and Levenshtein distance-based semantic parsing.

The system is designed to be backward compatible down to Android 7.0, and has undergone rigorous compatibility and performance testing on mainstream physical devices running Android 14 to 16. The testing demonstrates 60 fps smooth rendering on low-end phones, a WeChat one-tap dial success rate of 99.1\%, a network recovery guidance rate of 100\%, and memory usage consistently between 95 and 115MB. This work provides a complete engineering solution for age-friendly communication terminal development with notable practical and social value.